% Ví dụ một bảng đặc tả use case dùng macro trong config/usecase.tex
\begin{ucspec}{UC-01}{Đăng nhập}
\ucrow{Mã UC}{UC-01}
\ucrow{Tên UC}{Đăng nhập}
\ucrow{Mô tả}{Là người dùng, tôi muốn đăng nhập vào hệ thống để sử dụng các chức năng dành cho thành viên.}
\ucrow{Actor}{Guest}
\ucrow{Pre-Condition}{Người dùng chưa đăng nhập.}
\ucrow{Post-Condition}{Hệ thống tạo phiên đăng nhập.}
\ucrow{Basic Flow}{%
\begin{ucsteps}
\st{1.} Người dùng chọn Đăng nhập.
\st{2.} Hệ thống hiển thị biểu mẫu.
\st{3.} Người dùng nhập thông tin và xác nhận.
\st{4.} Hệ thống xác thực và tạo phiên đăng nhập.
\end{ucsteps}}
\ucrow{Alternative Flow}{%
\begin{ucsteps}
\st{4a.} Thông tin sai: hệ thống báo lỗi và quay lại bước 2.
\end{ucsteps}}
\ucrow{Business Rule}{%
\begin{ucsteps}
\st{BR1.} Quy tắc nghiệp vụ \chuachot{ví dụ điểm chưa chốt}.
\end{ucsteps}}
\end{ucspec}

% Chèn bộ ba sơ đồ activity/class/sequence (cần có ảnh trong thư mục diagrams):
% \ucdiagrams{uc-01}{UC-01}{Đăng nhập}
